\clearpage\begingroup
\providecommand{\ObsSec}[2]{#1\mathord{\overset{\mathrm{s}}{,}}#2}
\providecommand{\ObsArc}[2]{#1\mathord{\overset{\prime\prime}{,}}#2}
\providecommand{\ObsJupiter}{\mathord{\tikz[baseline=-.1ex,x=1em,y=1em,line width=.045em]{
  \draw (.02,.48) .. controls (.04,.70) and (.33,.69) .. (.31,.46) .. controls (.29,.29) and (.14,.19) .. (.02,.13)--(.49,.13);
  \draw (.37,.64)--(.37,-.06);}}}
\providecommand{\ObsVesta}{\mathord{\tikz[baseline=-.1ex,x=1em,y=1em,line width=.045em]{
  \draw (.04,.02) rectangle (.40,.38);
  \draw (.22,.38)--(.22,.48) (.12,.63)--(.22,.48)--(.32,.63);}}}

\SourceLeaf{292}\HistoricalHead{287}{R}{ÜBER DAS KREISMIKROMETER.}
\begin{center}\textsc{Gauss} an \textsc{Olbers}, Göttingen, 11. Februar 1808.\end{center}
\noindent\mbox{. . . . . . . . . . . . . . .}\par
\noindent\hspace*{\parindent}\textsc{Bessels} Methode\textsuperscript{[*)]}, Kreismikrometerbeobachtungen mit Rücksicht auf Refraction zu reduciren, kenne ich nicht, und kann also in Rücksicht auf die Schärfe zwischen derselben und der meinigen [keine Vergleichung] anstellen. Bei der meinigen liegen zwei nur näherungsweise richtige Voraussetzungen zum Grunde, [1.] dass das kreisförmige Gesichtsfeld durch die Refraction elliptisch wird (wenn ich mir diese \textls[120]{Phrase} erlauben darf), 2. dass das Verhältnis der Dimensionen der Ellipse zu denen des Kreises unabhängig von der Grösse des Gesichtsfeldes ist. Der letzten kann man überhoben sein, wenn man sich für sein Fernrohr eine besondere Tabelle construirt, wie die meinige, die eigentlich für ein unendlich kleines Feld gilt; der Unterschied wird aber bei einem nicht zu grossen Gesichtsfelde $=0$ sein. Ueber die Unrichtigkeit der ersten Voraussetzung will ich nächstens einmal eine Untersuchung anstellen; ich zweifle aber, dass dieselbe, vielleicht ganz kleine Höhen ausgenommen, nur so viel Einfluss haben könne, als die Behandlung der Kreismikrometerbeobb. überhaupt, insofern man sich erlaubt, den Parallel im Gesichtsfelde als geradlinig zu betrachten.\quad\mbox{. . . . . .}
\begingroup\renewcommand{\thefootnote}{[*)}\footnotetext{Monatl. Corresp, Bd. XVII, März 1808 ]}\endgroup
\begin{center}\textsc{Gauss} an \textsc{Olbers}, Göttingen, 13. Februar 1821.\end{center}
\noindent\mbox{. . . . . . . . . . . . . . .}\par
\begingroup\raggedright \noindent\hspace*{\parindent}Erlauben Sie mir noch ein Wort über die Berechnung der Kreismikrom[eter]beobachtungen nach dem Princip der Wahrscheinlichkeitsrechnung. Ist $\dfrac{15\,\mathrm{mora}.\cos\delta}{2\,\mathrm{rad.}}=\cos\varphi$ für den Stern u[nd] $\varphi'$ dasselbe für den Kometen, so ist Decl.-Unt[erschied] $=r\sin\varphi\mp r\sin\varphi'$. Das dem Rectascensionsunterschiede beizulegende Gewicht ist dann $\dfrac{2}{\dfrac{1}{\cos\varphi^2}+\dfrac{1}{\cos\varphi^{\prime 2}}}$ u[nd] das Gewicht der Decl.-Best[immung] $=\dfrac{2}{\dfrac{1}{\sin\varphi^2}+\dfrac{1}{\sin\varphi^{\prime 2}}}$.\par\endgroup

Hiebei ist freilich nur der Fehler des Sehens berücksichtigt und die Beob. des $\ast$ u[nd] Kometen als gleich leicht angesehen. (Man könnte die Verschiedenheit leicht berücksichtigen, wenn sie in Zahlen bekannt wäre.)
\clearpage\SourceLeaf{293}\HistoricalHead{288}{L}{PRAKTISCHE UND SPHÄRISCHE ASTRONOMIE. BRIEFWECHSEL.}
\noindent Auch das Gesichtsfeld ist als genau bekannt betrachtet. Aber auch ohne diese Umstände zu berücksichtigen, wird es doch viel genauer sein, die Beobachtungen nach jenen Gewichten zu berechnen, als wenn man alle Resultate als gleich genau ansieht. Ich habe dies bei den Cometenbeobb. gethan. Die Rechnung ist höchst einfach, wenn man sich meiner kleinen Logarithmentafeln bedient; ich gehe für die erste Formel mit der Differenz der Logarithmen von $\cos\varphi^2$ u[nd] $\cos\varphi^{\prime 2}$ in die Kolumne $A$ ein und subtrahire $B$ allemal von dem \textls[120]{kleinern} Logarithmen von $2\cos\varphi^2$ oder $2\cos\varphi^{\prime 2}$. Ebenso bei der andern Formel. Am 10. Februar waren z. B. meine einzelnen Resultate\textsuperscript{[*)]}
\begin{center}\renewcommand{\arraystretch}{1.15}\begin{tabular}{rlr|rlr}
\hline\hline
\multicolumn{3}{c|}{für Decl.-Unterschied}&\multicolumn{3}{c}{für $AR$}\\
\hline\hline
$\ObsArc{135}{5}$&Gewicht&$0{,}666$&$\ObsArc{47}{25}$&Gewicht&$0{,}294$\\
$126{,}0$&»&$0{,}889$&$52{,}05$&»&$0{,}001$\\
$140{,}7$&»&$0{,}138$&$47{,}95$&»&$1{,}000$\\
$149{,}9$&»&$0{,}614$&$47{,}90$&»&$0{,}812$\\
$134{,}1$&»&$0{,}719$&$48{,}00$&»&$0{,}240$\\
&&&$48{,}25$&»&$1{,}000$\\
&&&$47{,}55$&»&$1{,}000$\\
&&&$48{,}25$&»&$0{,}116$
\end{tabular}\end{center}
\noindent\hspace*{\parindent}Also die Gewichte \textls[120]{sehr} ungleich. Bei der Decl. müssen die sehr nahe beim Centrum geschehenen Durchgänge ganz verworfen werden, aber von den $AR$ braucht man gar keine auszuschliessen.

Das Gewicht einer $AR$, die für $\ast$ u[nd] Komet central ist, ist hiebei als Einheit angenommen, sowie das einer Decl.-Bestimmung, wo $\ast$ und Komet hart am Rande hinstreichen.
\begingroup\renewcommand{\thefootnote}{[*)}\footnotetext{Beobachtungen des Kometen 1821, Berl. Astron. Jahrbuch für 1825, Werke VI, S. 439.]}\endgroup
\begin{center}BEMERKUNGEN.\end{center}
\begingroup\small
\noindent\hspace*{\parindent}Zum Brief vom 5. Januar 1808. — Setzt man voraus, dass das Kreismikrometer aus der Himmelskugel eine Ellipse ausschneidet, die als eben betrachtet werden kann, so ist eine Beziehung zwischen der Länge der vom Stern beschriebenen Sehne und ihrem Abstand von der Mitte der Ellipse herzustellen. Man bezeichne die erstere mit $\mu$ und den letzteren mit $d$, lege ein rechtwinkliges Koordinatensystem durch die beiden Axen $a$ und $b$ der Ellipse und nenne $\alpha$ den Winkel, unter dem die Sehne die $x$-Axe schneidet.
\clearpage\SourceLeaf{294}\HistoricalHead{289}{R}{ÜBER DAS KREISMIKROMETER. BEMERKUNGEN.}
\noindent Dann sind die Gleichungen der Ellipse und der Sehne
\[x^2(1-\lambda^2)+y^2=b^2,\qquad y\cos\alpha=x\sin\alpha+d,\]
\noindent also die Schnittpunkte beider entsprechend dem Ein- und Austritt des Sterns
\[x(1-\lambda^2\cos\alpha^2)=-d\sin\alpha\pm\surd\bigl\{(1-\lambda^2)(a^2(1-\lambda^2\cos\alpha^2)-d^2)\bigr\}\]
\noindent und hieraus folgt für die Länge der Sehne
\[\frac14\mu^2=\frac14\left(\frac{x_2-x_1}{\cos\alpha}\right)^2=\frac{(1-\lambda^2)(a^2(1-\lambda^2\cos\alpha^2)-d^2)}{(1-\lambda^2\cos\alpha^2)^2}.\]
\noindent\hspace*{\parindent}Führt man den Winkel $\psi$ ein, indem man setzt
\[d^2=a^2\sin\psi^2(1-\lambda^2\cos\alpha^2),\]
\noindent so wird
\[\frac14\mu^2=\frac{a^2(1-\lambda^2)\cos\psi^2}{1-\lambda^2\cos\alpha^2}.\]
\noindent\hspace*{\parindent}Da
\[d=\delta-D,\qquad\mu=(t'-t)\cos\delta,\qquad p=90^0-\alpha,\]
\noindent so folgen die beiden \textsc{Gauss}schen Gleichungen I und II aus den beiden vorstehenden. Ferner hat man für die Abszisse der Sehnenmitte
\[\frac12(x_2+x_1)=-\frac{d\sin\alpha}{1-\lambda^2\cos\alpha^2}.\]
\noindent Die Abszisse des Fusspunktes des vom Mittelpunkt auf die Sehne gefällten Lotes ist gleich
\[-d\sin\alpha.\]
\noindent Der Abstand beider also $=\dfrac{d\lambda^2\sin\alpha\cos\alpha}{1-\lambda^2\cos\alpha^2}$ entsprechend der \textsc{Gauss}schen Gleichung III.

Als Kreismikrometer diente seinerzeit der Rand des Gesichtsfeldes.

Zum Brief vom 13. Februar 1821. — Bedeutet $\varphi$ den Winkel zwischen dem Parallel und dem nach dem Eintrittsort gezogenen Radius des Gesichtsfeldes, so ist die Eintrittsdauer, und damit der mittlere Fehler der Eintrittszeit, $=\dfrac{1}{\cos\varphi}$, wenn dieser bei senkrechtem Eintritt $=1$ gesetzt wird. Hiernach wird das mittlere Fehlerquadrat der AR-Differenz von Stern und Objekt $=\dfrac{1}{\cos\varphi^2}+\dfrac{1}{\cos\varphi^{\prime 2}}$, also das Gewicht wie im obigen Briefe. In gleicher Weise folgt das Gewicht der Deklinationsbestimmung. Auf die Deklination — den Faktor $\cos\delta$ — ist dabei keine Rücksicht genommen.

Mit dem Ausdruck »mora« bezeichnet \textsc{Gauss} die Durchgangsdauer $t'-t$.
\begin{flushright}\textsc{Brendel}.\end{flushright}
\begin{center}\rule{.31\textwidth}{.35pt}\end{center}
\clearpage\endgroup\endgroup
